% !TEX root = ../main.tex
\subsection{Training process}

The questions were split at random (seed 24) into training (12,909), validation (1,614) and test (1,614)
sets. The training set is the catalog; the others act as user queries.

No human relevance labels exist, so a returned question counts as relevant if it shares an AI/ML tag with
the query. We report precision at $k$ (P@k), hit rate at $k$ (Hit@k: at least one relevant result in the
top $k$), mean reciprocal rank (MRR) and the precision of random ranking (0.073).

The TF-IDF vectorizer was fitted on the training set only, so nothing leaked from the evaluation sets.
Four configurations were compared on validation P@5 (Table~\ref{tab:grid}) and the best was kept. The
scores barely differ (0.604 to 0.610), so the method, not its settings, limits TF-IDF.

\begin{table}[h]
  \centering\small
  \caption{TF-IDF configurations compared on validation P@5. All use character 3--5-grams; the last row was selected.}
  \label{tab:grid}
  \begin{tabular}{cccc c}
    \toprule
    Word n-grams & min\_df & max\_features & Stop-words & Val.\ P@5 \\
    \midrule
    1--1 & 2 & 100k & none    & 0.609 \\
    1--2 & 2 & 100k & none    & 0.606 \\
    1--2 & 3 & 50k  & none    & 0.604 \\
    1--2 & 2 & 100k & English & 0.610 \\
    \bottomrule
  \end{tabular}
\end{table}

MiniLM was used pre-trained: we have no labelled pairs of matching questions, and shared tags are too
coarse to fine-tune on. Its training step is indexing, where each catalog question is encoded once and
the vectors are saved for the app (Listing~\ref{lst:encode}).

\begin{lstlisting}[caption={Encoding the catalog with MiniLM and scoring test queries (\texttt{train\_model.ipynb}).}, label={lst:encode}]
embedding_model = SentenceTransformer("all-MiniLM-L6-v2")
train_embeddings = embedding_model.encode(train_df["text_all"].tolist(), convert_to_numpy=True)
similarity = cosine_similarity(test_embeddings, train_embeddings)   # 500 test queries x 12,909
\end{lstlisting}
